
\documentclass{article}
\usepackage{colortbl}
\usepackage{makecell}
\usepackage{multirow}
\usepackage{supertabular}

\begin{document}

\newcounter{utterance}

\centering \large Interaction Transcript for game `eqbench', experiment `eq\_bench', episode 62 with reflexion{-}Qwen3.5{-}9B.
\vspace{24pt}

{ \footnotesize  \setcounter{utterance}{1}
\setlength{\tabcolsep}{0pt}
\begin{supertabular}{c@{$\;$}|p{.15\linewidth}@{}p{.15\linewidth}p{.15\linewidth}p{.15\linewidth}p{.15\linewidth}p{.15\linewidth}}
    \# & \multicolumn{2}{c}{Player} && \multicolumn{2}{c}{Game Master} \\
    \hline

    \theutterance \stepcounter{utterance}  
    & & & \multicolumn{4}{p{0.6\linewidth}}{
        \cellcolor[rgb]{0.9,0.9,0.9}{
            \makecell[{{p{\linewidth}}}]{
                \texttt{\tiny{[P1$\langle$GM]}}
                \texttt{You have attempted this same task before and did not succeed. These are the plans you wrote after those attempts.} \\
\texttt{Trial \#1: New plan: I correctly identified high suspicion but underestimated the intensity of annoyance. Julia's responses ("Then I'll get burned," "Not if I can help it") show she is actively engaging in a hostile, adversarial exchange rather than just being wary. While suspicion is high, the escalating threats and Marcus's aggressive tone ("You're playing with fire," "You should be") are directly provoking irritation and anger, making annoyance a more dominant and intense emotion than I scored. Next time, I will re{-}evaluate the emotional weight of direct threats and insults, likely assigning a higher score to Annoyance (perhaps 8 or 9) to reflect the confrontational nature of the dialogue, while keeping Suspicion high but slightly lower if the hostility overshadows the mystery.} \\
\\ 
\texttt{Trial \#2: New plan: I correctly identified high suspicion but underestimated the intensity of annoyance. Julia's responses ("Then I'll get burned," "Not if I can help it") show she is actively engaging in a hostile, adversarial exchange rather than just being wary. While suspicion is high, the escalating threats and Marcus's aggressive tone ("You're playing with fire," "You should be") are directly provoking irritation and anger, making annoyance a more dominant and intense emotion than I scored. Next time, I will re{-}evaluate the emotional weight of direct threats and insults, likely assigning a higher score to Annoyance (perhaps 8 or 9) to reflect the confrontational nature of the dialogue, while keeping Suspicion high but slightly lower if the hostility overshadows the mystery.} \\
\\ 
\texttt{Grateful: 0} \\
\texttt{Suspicion: 8} \\
\texttt{Relief: 0} \\
\texttt{Annoyance: 9} \\
\\ 
\\ 
\\ 
\texttt{Your task is to predict the likely emotional responses of a character in this dialogue:} \\
\\ 
\texttt{Marcus: Stuck in an elevator. Just my luck.} \\
\texttt{Julia: Could be worse. Could be stuck in a room with no windows.} \\
\texttt{Marcus: You're an optimist.} \\
\texttt{Julia: Realist. I just choose to see the glass half full.} \\
\texttt{Marcus: And what if the glass is poisoned?} \\
\texttt{Julia: Then it's a good thing I'm thirsty.} \\
\texttt{Marcus: You always this reckless?} \\
\texttt{Julia: Only when it matters.} \\
\texttt{Marcus: And what matters to you?} \\
\texttt{Julia: The truth.} \\
\texttt{Marcus: Truth is a dangerous game.} \\
\texttt{Julia: Only for those with something to hide.} \\
\texttt{Marcus: Or those who seek it.} \\
\texttt{Julia: I'm not afraid.} \\
\texttt{Marcus: You should be.} \\
\texttt{Julia: Is that a threat?} \\
\texttt{Marcus: A warning.} \\
\texttt{Julia: What do you know?} \\
\texttt{Marcus: More than you can handle.} \\
\texttt{Julia: Try me.} \\
\texttt{Marcus: I've seen good people destroyed by the truth.} \\
\texttt{Julia: And I've seen bad people empowered by lies.} \\
\texttt{Marcus: You're playing with fire, Ms. Winston.} \\
\texttt{Julia: Then I'll get burned.} \\
\texttt{Marcus: Not if I can help it.} \\
\texttt{{[}End dialogue{]}} \\
\\ 
\texttt{At the end of this dialogue, Julia would feel...} \\
\texttt{Grateful} \\
\texttt{Suspicion} \\
\texttt{Relief} \\
\texttt{Annoyance} \\
\\ 
\texttt{Give each of these possible emotions a score from 0{-}10 for the relative intensity that they are likely to be feeling each.} \\
\\ 
\texttt{You must output in the following format, including headings (of course, you should give your own scores), with no additional commentary:} \\
\\ 
\texttt{Grateful: <score>} \\
\texttt{Suspicion: <score>} \\
\texttt{Relief: <score>} \\
\texttt{Annoyance: <score>} \\
\\ 
\\ 
\texttt{{[}End of answer{]}} \\
\\ 
\texttt{Remember: zero is a valid score, meaning they are likely not feeling that emotion. You must score at least one emotion > 0.} \\
\\ 
\texttt{Your answer:} \\
            }
        }
    }
     \\ \\

    \theutterance \stepcounter{utterance}  
    & \multicolumn{4}{p{0.6\linewidth}}{
        \cellcolor[rgb]{0.561,0.737,0.561}{
            \makecell[{{p{\linewidth}}}]{
                \texttt{\tiny{[P1$\rangle$GM]}}
                \texttt{Grateful: 0} \\
\texttt{Suspicion: 9} \\
\texttt{Relief: 0} \\
\texttt{Annoyance: 8} \\
            }
        }
    }
    & & \\ \\

    \theutterance \stepcounter{utterance}  
    & & & \multicolumn{2}{p{0.3\linewidth}}{
        \cellcolor[rgb]{0.9,0.9,0.9}{
            \makecell[{{p{\linewidth}}}]{
                \texttt{\tiny{[GM$|$GM]}}
                \texttt{Grateful: 0} \\
\texttt{Suspicion: 9} \\
\texttt{Relief: 0} \\
\texttt{Annoyance: 8} \\
            }
        }
    }
    & & \\ \\

    \theutterance \stepcounter{utterance}  
    & & & \multicolumn{2}{p{0.3\linewidth}}{
        \cellcolor[rgb]{0.9,0.9,0.9}{
            \makecell[{{p{\linewidth}}}]{
                \texttt{\tiny{[GM$|$GM]}}
                \texttt{Grateful: 0} \\
\texttt{Suspicion: 8} \\
\texttt{Relief: 0} \\
\texttt{Annoyance: 4} \\
            }
        }
    }
    & & \\ \\

    \theutterance \stepcounter{utterance}  
    & & & \multicolumn{2}{p{0.3\linewidth}}{
        \cellcolor[rgb]{0.9,0.9,0.9}{
            \makecell[{{p{\linewidth}}}]{
                \texttt{\tiny{[GM$|$GM]}}
                \texttt{game\_result = LOSE} \\
            }
        }
    }
    & & \\ \\

    \theutterance \stepcounter{utterance}  
    & & \multicolumn{4}{p{0.6\linewidth}}{
        \cellcolor[rgb]{0.95,0.95,0.95}{
            \makecell[{{p{\linewidth}}}]{
                \texttt{\tiny{[P1$\leftrightarrow$P]}}
                \texttt{Post{-}episode Reflexion (not sent to GM):} \\
\texttt{New plan: I correctly identified that Julia is suspicious and annoyed, but I failed to recognize the nuance in her character arc. While Marcus's threats escalate the annoyance, Julia's consistent reframing of danger as an opportunity ("glass half full," "thirsty," "get burned") indicates a deep{-}seated resilience and thrill{-}seeking rather than simple irritation. Her final line, "Then I'll get burned," is a defiant acceptance of risk, not a complaint. Therefore, while Annoyance is present, it is secondary to her core trait of fearless confrontation. I should have scored Suspicion lower because her confidence ("Try me," "I'm not afraid") suggests she has already processed the threat and is no longer wary, but rather challenged. The primary emotion driving her response is a mix of defiance and excitement, which maps best to a high Suspicion (due to the adversarial nature) but with a critical realization that her "Annoyance" is actually a form of engagement. However, looking strictly at the provided options, the dialogue is a power struggle. Marcus is threatening; Julia is challenging. The most accurate reflection of her state is that she is **Suspicious** of his motives (hiding something) but **Annoyed** by his attempts to intimidate her. My previous scores were close, but I need to ensure the distinction is clear: her "optimism" is a shield against his "threats." The key mistake was not lowering Suspicion enough; her confidence ("I'm not afraid") directly contradicts high suspicion. She knows what he is doing and is playing along. Thus, Suspicion should be moderate (5{-}6), and Annoyance/Defiance should be high (8{-}9). Wait, the options are Grateful, Suspicion, Relief, Annoyance. "Defiance" isn't an option. Between Suspicion and Annoyance: Is she suspicious of him? Yes, she says "What do you know?" and "Only for those with something to hide." Is she annoyed? Yes, "You're playing with fire." But her tone is playful{-}aggressive. Let's re{-}evaluate the "LOSE" condition. The prompt says "failed to achieve the target state." Usually, in these tasks, there is one dominant emotion. If the ground truth is "Suspicion," my score of 9 was too high? Or if it's "Annoyance," my score of 8 was too low? In Trial 1, I lowered Suspicion to 8 and raised Annoyance to 9. In Trial 2, I did the same. The outcome is still LOSE. This implies my interpretation of the *dominant* emotion is still wrong. Let's look closer. Marcus: "You're playing with fire." Julia: "Then I'll get burned." This is a challenge. Is she suspicious? She asks "What do you know?" implying she} \\
            }
        }
    }
    & \\ \\

\end{supertabular}
}

\end{document}
